\MRPMAPGraph{}

\MRPLMERGraph{}

\MRPLMERSdGraph{}

In \figref{mrp_map_graph} we see that the MAP estimator provides poor point
estimates in the MrP model, likely due in part to the large number of random
effects to fit.  We note that the \texttt{stan::optimizing} function in \texttt{R}
performed poorly, and we had find the MAP by hand using a careful initialization
and small step size in order to avoid divergent estimates.

Although it is not strictly speaking a MAP estimate, the natural
optimization--based competitor for MCMC is not the MAP estimator but
marginal maximum likelihood estimation, such as the \texttt{lmer} function
in \texttt{R}.  As shown in \figref{mrp_lmer_graph}, \texttt{lmer} provides
better point estimates than MAP estimation on the MrP model.  However,
as shown in \figref{mrp_lmer_sd_graph}, the standard error estimates are
provided by \texttt{lmer} can differ considerably from the Bayesian posterior
estimates, and most parameters do not have standard errors reported automatically.
In principle, one may be able combine the frequentist sampling estimates of each
parameter to estimate the joint sampling distribution between the fixed
in order to estimate the sampling variance of $\expect{\post}{\g(\theta)}$,
but doing so would be quite cumbersome given the information exposed by \texttt{lmer}.
In contrast, the MCMC and IJ covariance estimates are extremely straightforward
to compute. 